\chapter{ALU 与标志位}

ALU 是算术逻辑单元。它不是凭空出现的“大部件”，而是因为加法、减法、按位逻辑、移位和比较反复出现，硬件需要把这些能力集中到一个受控制的组合逻辑块中。ALU 的本质是：多个候选计算电路并排存在，控制信号选择本周期要看的结果。

\section{一、ALU 的输入输出先从需求看}

假设一块硬件已经能保存两个 8-bit 值 A 和 B。它可能需要 A+B，也可能需要 A-B，也可能只需要 A AND B 来屏蔽某些 bit，还可能需要判断 A 是否等于 B。若每个需求都单独拉一条完整路径，后面的连线和控制会很乱。

更自然的做法是给这组计算一个统一入口：

\begin{lstlisting}
ALU(a, b, op) -> result, flags
\end{lstlisting}

a 和 b 是两个操作数，op 是选择信号，result 是结果，flags 是结果的附带证据。一个最小 ALU 可以包含：

\begin{longtable}{p{0.18\linewidth}p{0.32\linewidth}p{0.38\linewidth}}
\toprule
op & 运算 & 硬件来源 \\
\midrule
ADD & a + b & 加法器 \\
SUB & a - b & 加法器加反相控制 \\
AND & a \& b & 按位 AND 门阵列 \\
OR & a | b & 按位 OR 门阵列 \\
XOR & a xor b & 按位 XOR 门阵列 \\
SHIFT & 移位 & 移位器 \\
\bottomrule
\end{longtable}

\begin{pdfdiagram}[x=1cm,y=1cm]
  \node[hw state,minimum width=1.1cm] (a) at (0.8,2.1) {a};
  \node[hw state,minimum width=1.1cm] (b) at (0.8,0.9) {b};
  \node[hw control,minimum width=1.0cm] (op) at (0.8,-0.2) {op};
  \node[hw compute,minimum width=1.45cm] (add) at (3.0,2.65) {ADD/SUB};
  \node[hw compute,minimum width=1.45cm] (logic) at (3.0,1.75) {AND/OR/XOR};
  \node[hw compute,minimum width=1.45cm] (shift) at (3.0,0.85) {SHIFT};
  \node[hw compute,minimum width=1.3cm] (cmp) at (3.0,-0.05) {比较证据};
  \node[hw compute,minimum width=1.4cm] (mux) at (5.6,1.35) {结果 mux};
  \node[hw state,minimum width=1.2cm] (result) at (7.9,1.35) {result};
  \node[hw state,minimum width=1.2cm] (flags) at (7.9,0.15) {flags};
  \foreach \target in {add,logic,shift,cmp} {
    \draw[hw bus] (a.east) -- ++(0.6,0) |- (\target.west);
    \draw[hw bus] (b.east) -- ++(0.85,0) |- (\target.west);
  }
  \draw[hw bus] (op.east) -- ++(2.7,0) |- (mux.south);
  \draw[hw bus] (add.east) -- (mux.west);
  \draw[hw bus] (logic.east) -- (mux.west);
  \draw[hw bus] (shift.east) -- (mux.west);
  \draw[hw bus] (cmp.east) -- (mux.west);
  \draw[hw bus] (mux.east) -- (result.west);
  \draw[hw bus] (add.east) .. controls (5.3,0.1) .. (flags.west);
  \draw[hw bus] (result.south) -- (flags.north);
\end{pdfdiagram}
\diagramnote{ALU 把多种候选结果并排产生，再由 op 选择最终结果；标志位来自结果和关键内部证据。}

\section{二、减法复用加法器}

ALU 里最值得细看的是 SUB。若单独做减法器，会增加面积和连线。补码让它复用 ADD 的加法器：

\begin{lstlisting}
if op = ADD:
  b_to_adder = b
  carry_in = 0
if op = SUB:
  b_to_adder = NOT b
  carry_in = 1
result = a + b_to_adder + carry_in
\end{lstlisting}

这段伪代码描述的是硬件选择。实际电路里，b 的每一位前面可以接 XOR 门：当 sub 控制为 0，b 原样通过；当 sub 控制为 1，b 被取反。最低位进位输入也由 sub 控制。这样 ADD 和 SUB 共用同一条进位链。

这个案例很重要，因为它展示了硬件设计的常见思路：先找到数学规则之间的共用部分，再用少量控制信号复用同一块组合逻辑。

\section{三、结果选择来自多路选择器}

ALU 内部可以并行产生多个候选结果，再用多路选择器按 op 选出最终 result。

\begin{lstlisting}
add_result
and_result
or_result
shift_result
result = mux(op, add_result, and_result, or_result, shift_result)
\end{lstlisting}

这不表示所有运算都“按顺序算一遍”。组合逻辑是并行传播的：加法器、按位逻辑、移位器都可能同时产生自己的候选输出。最终哪一路被后级看见，由 mux 决定。

代价也随之出现。候选电路越多，ALU 面积越大；mux 输入越多，选择延迟越大；若某个候选结果来自长进位链，它可能成为 ALU 的关键路径。因此 ALU 不是“把功能塞进去就完了”，还要考虑延迟、面积和功耗。

\section{四、标志位是结果证据}

ALU 常见标志位包括：

\begin{longtable}{p{0.2\linewidth}p{0.38\linewidth}p{0.3\linewidth}}
\toprule
标志位 & 含义 & 用途 \\
\midrule
Zero & result 是否为 0 & 相等判断、条件选择 \\
Carry & 无符号进位或借位证据 & 多字长加法、无符号比较 \\
Overflow & 补码有符号溢出 & 有符号范围检查 \\
Negative & result 最高位是否为 1 & 补码符号证据 \\
\bottomrule
\end{longtable}

Zero 可以由所有结果 bit 做 OR 后再取反得到。只要有任意一位为 1，结果就不是 0；所有位都是 0，Zero 才为 1。Negative 通常直接来自结果最高位。Carry 来自加法器最高位外的进位。Overflow 则要看有符号加减法中符号是否出现不可能的变化。

标志位不是额外装饰，而是后续控制的重要输入。若比较 A 和 B 是否相等，不一定要把完整结果拿出去再分析；ALU 做 A-B 后，Zero=1 就说明两者相等。

例如要判断两个 8-bit 值是否相等，可以让 ALU 执行减法。若结果 8 个 bit 全部为 0，Zero 置 1。控制器看到 Zero=1，就知道两个输入相等。这个判断没有另外发明一套“相等机器”，而是复用了减法路径和结果检测网络。

\section{五、ALU 是组合部件，保存要靠寄存器}

ALU 本身通常是组合逻辑。输入一变，结果经过传播延迟后改变。它不会记住上一次运算。若要保存 ALU 结果，必须在输出后接寄存器。

这个边界必须分清：ALU 负责计算，寄存器负责保存，控制信号负责选择。下一章会把 ALU 接到寄存器阵列、选择器和内部连线上，形成数据通路。没有数据通路，ALU 只是孤立计算块；没有 ALU，数据通路只能搬运，不能变换。

\begin{classicnote}
ALU 的经典读法是“候选计算并行产生，选择器选出一个结果，标志位保存结果证据”。不要把它想成一段顺序执行的程序。

\begin{itemize}
\item 纸上实验：给定 \code{alu\_op} 四种取值 ADD、SUB、AND、XOR，写出每种情况下结果、Zero、Negative、Carry、Overflow 的来源。
\item 机器证据：减法复用加法器时，B 输入前的反相器和 carry-in=1 是控制信号驱动的真实硬件路径。
\item 控制证据：分支不一定需要完整结果写回寄存器，控制器可能只读取 Zero 或 Negative 等标志来选择下一 PC。
\item 检查题：若 ALU 结果已经正确，但 Zero 标志生成逻辑漏接最高位，哪类指令最先暴露问题？
\end{itemize}

经典教材会把 ALU 看成数据通路的“可配置组合核心”。它不是 CPU 全部，只是一个在控制线约束下把输入 bit 变成输出 bit 和证据 bit 的部件。
\end{classicnote}
